\begin{table}[htbp]
\centering
\footnotesize
\caption{Problema G, Passos 6 e 7 (openai/gpt-4o-mini): LLM versus Perceptron treinado nos mesmos exemplos, acurácia contra o rótulo real.}
\label{tab:bloco3_G_passo7_LLM_vs_perceptron_gpt4mini}
\begin{tabular}{llllrrrrrrrrrr}
\toprule
 &  &  &  & \multicolumn{5}{c}{Acur.\ LLM vs real (\%)} & \multicolumn{5}{c}{Acur.\ Perceptron vs real (\%)} \\ \cmidrule(lr){5-9} \cmidrule(lr){10-14}
Subproblema & Classes & Atributos & Prompt & k = 0 & k = 4 & k = 10 & k = 20 & k = 40 & k = 0 & k = 4 & k = 10 & k = 20 & k = 40 \\
\midrule
G1 (normalizado) & Mulher / Homem & peso/altura & R$^2$ & 43,3 & 79,6 & 84,8 & 86,7 & 88,9 & — & 79,3 & 80,0 & 81,5 & 83,0 \\
G1 (normalizado) & Mulher / Homem & peso/altura & R$^3$ & 43,3 & 70,7 & 84,1 & 89,3 & 91,1 & — & 78,1 & 80,4 & 81,5 & 82,6 \\
G1 (normalizado) & Mulher / Homem & peso/altura & R$^4$ & 43,3 & 73,0 & 77,8 & 80,0 & 85,9 & — & 78,1 & 80,0 & 80,7 & 82,2 \\
G1 (normalizado) & Mulher / Homem & x1/x2 & R$^2$ & 43,3 & 76,3 & 84,1 & 85,6 & 88,9 & — & 79,3 & 80,0 & 81,5 & 83,0 \\
G1 (normalizado) & Mulher / Homem & x1/x2 & R$^3$ & 43,3 & 73,7 & 86,3 & 90,0 & 91,1 & — & 78,1 & 80,4 & 81,5 & 82,6 \\
G1 (normalizado) & Mulher / Homem & x1/x2 & R$^4$ & 43,3 & 75,2 & 79,6 & 83,3 & 89,3 & — & 78,1 & 80,0 & 80,7 & 82,2 \\
G1 (normalizado) & A / B & peso/altura & R$^2$ & 24,4 & 81,1 & 87,4 & 87,8 & 88,5 & — & 80,0 & 81,1 & 83,0 & 84,1 \\
G1 (normalizado) & A / B & peso/altura & R$^3$ & 24,4 & 70,7 & 86,7 & 87,8 & 88,5 & — & 78,9 & 81,5 & 83,0 & 84,8 \\
G1 (normalizado) & A / B & peso/altura & R$^4$ & 24,4 & 74,8 & 82,6 & 84,4 & 88,1 & — & 79,3 & 81,1 & 81,9 & 84,1 \\
G1 (normalizado) & A / B & x1/x2 & R$^2$ & 31,1 & 80,4 & 86,3 & 87,4 & 88,5 & — & 79,3 & 79,6 & 80,4 & 82,6 \\
G1 (normalizado) & A / B & x1/x2 & R$^3$ & 31,1 & 73,3 & 85,2 & 88,5 & 89,3 & — & 78,1 & 80,0 & 80,4 & 82,2 \\
G1 (normalizado) & A / B & x1/x2 & R$^4$ & 31,1 & 76,3 & 82,6 & 85,2 & 89,6 & — & 78,1 & 79,6 & 79,3 & 81,5 \\
\textbf{G2 (kg e m)} & \textbf{Mulher / Homem} & \textbf{peso\_kg/altura\_m} & \textbf{R$^2$} & \textbf{91,1} & \textbf{80,7} & \textbf{81,1} & \textbf{74,1} & \textbf{77,8} & \textbf{—} & \textbf{79,3} & \textbf{80,0} & \textbf{81,1} & \textbf{82,6} \\
\textbf{G2 (kg e m)} & \textbf{Mulher / Homem} & \textbf{peso\_kg/altura\_m} & \textbf{R$^3$} & \textbf{91,1} & \textbf{74,1} & \textbf{75,2} & \textbf{74,1} & \textbf{75,2} & \textbf{—} & \textbf{78,1} & \textbf{80,4} & \textbf{81,5} & \textbf{82,2} \\
\textbf{G2 (kg e m)} & \textbf{Mulher / Homem} & \textbf{peso\_kg/altura\_m} & \textbf{R$^4$} & \textbf{91,1} & \textbf{76,7} & \textbf{75,9} & \textbf{73,0} & \textbf{78,5} & \textbf{—} & \textbf{78,1} & \textbf{80,0} & \textbf{80,4} & \textbf{82,6} \\
G2 (kg e m) & Mulher / Homem & x1/x2 & R$^2$ & 63,3 & 73,3 & 73,7 & 70,7 & 74,4 & — & 79,3 & 80,0 & 81,5 & 83,0 \\
G2 (kg e m) & Mulher / Homem & x1/x2 & R$^3$ & 63,3 & 71,5 & 73,7 & 73,0 & 74,4 & — & 78,1 & 80,4 & 81,5 & 82,6 \\
G2 (kg e m) & Mulher / Homem & x1/x2 & R$^4$ & 63,3 & 72,2 & 72,2 & 77,8 & 80,7 & — & 78,1 & 80,0 & 80,7 & 82,2 \\
G2 (kg e m) & A / B & peso\_kg/altura\_m & R$^2$ & 40,0 & 74,8 & 76,7 & 74,1 & 76,3 & — & 79,3 & 80,0 & 81,1 & 82,6 \\
G2 (kg e m) & A / B & peso\_kg/altura\_m & R$^3$ & 40,0 & 72,6 & 75,2 & 76,3 & 76,7 & — & 78,1 & 80,4 & 81,5 & 82,6 \\
G2 (kg e m) & A / B & peso\_kg/altura\_m & R$^4$ & 40,0 & 74,1 & 76,3 & 79,3 & 84,8 & — & 78,1 & 80,0 & 80,4 & 83,0 \\
G2 (kg e m) & A / B & x1/x2 & R$^2$ & 43,3 & 70,7 & 74,1 & 75,6 & 77,4 & — & 80,0 & 81,1 & 81,9 & 83,0 \\
G2 (kg e m) & A / B & x1/x2 & R$^3$ & 43,3 & 70,4 & 77,4 & 77,4 & 78,9 & — & 78,9 & 81,5 & 81,9 & 83,0 \\
G2 (kg e m) & A / B & x1/x2 & R$^4$ & 43,3 & 70,7 & 75,6 & 83,0 & 85,6 & — & 79,3 & 81,1 & 80,7 & 83,0 \\
\bottomrule
\multicolumn{14}{p{0.95\linewidth}}{\footnotesize Problema G, Passo 6, nas 30 amostras de teste; o especialista é a métrica eleita no Passo 5. Prompt = atributos mostrados ao LLM (R$^2$ = x1, x2; R$^3$ = + x1·x2; R$^4$ = + x1², x2²), fator controlado independente da métrica eleita. Perceptron (Passo 7) = Perceptron diagonal treinado nos mesmos $k$ exemplos, no mesmo espaço que o LLM viu (coluna Prompt); se todos os regimes dos Passos 2 e 3 colapsaram, sem métrica eleita, o Passo 6 não roda e a linha não existe. Acurácia (vs rótulo real) do LLM e do Perceptron; médias das sementes e sorteios. Nas linhas com classes neutras A/B a acurácia literal do LLM pode sair invertida (ver bloco3\_table\_02). Com 3 ou 4 atributos no prompt, os exemplos do G2 mostram os extras em unidades físicas (ex.: peso\_kg·altura\_m), que não são função monótona de x1·x2 normalizado.} \\
\end{tabular}
\end{table}
